\documentclass[11pt]{article}
\usepackage[margin=1in]{geometry}
\usepackage{amsmath}
\usepackage{booktabs}
\usepackage{hyperref}

\title{Retrieval-Grounded Assistance for Reproducible Literature Reviews}
\author{ScholarAI Example Author}
\date{}

\begin{document}
\maketitle

\begin{abstract}
Large language models can accelerate literature reviews, but fluent output alone does not guarantee that a claim is supported by a source. This short sample describes a retrieval-grounded workflow that keeps evidence, drafting, and human review connected. It is intentionally compact so that users can run the complete ScholarAI paper workflow and inspect every input and output.
\end{abstract}

\section{Introduction}
Systematic reviews require transparent search and screening procedures. The PRISMA 2020 statement provides updated reporting guidance for this process \cite{page2021prisma}. Recent retrieval-augmented generation systems offer another useful building block: they condition generated text on retrieved documents instead of relying only on model parameters \cite{lewis2020retrieval}. However, a practical research assistant must also expose where evidence came from and where human confirmation is still required.

\section{Method}
Our example workflow has three stages. First, a retriever returns a candidate set $D_q$ for a research question $q$. Second, each candidate receives a relevance score $s(d,q)$. Finally, the drafting model receives only the top-ranked evidence:
\begin{equation}
  D_q^{(k)} = \operatorname{TopK}_{d \in D_q} s(d,q).
\end{equation}
Every generated recommendation stores its supporting document identifiers. Recommendations without traceable evidence are marked for manual review rather than silently merged into the manuscript.

\section{Evaluation Plan}
The workflow can be evaluated on retrieval recall, citation correctness, and reviewer effort. A useful experiment compares an ungrounded baseline with the retrieval-grounded workflow while keeping the language model fixed. Reviewers then label whether each factual statement is supported, contradicted, or not addressed by its cited source.

\section{Conclusion}
Retrieval grounding does not remove the need for expert judgment, but it makes the boundary between evidence and generation easier to inspect. This sample is a starting point: users should replace its text and references with material appropriate to their own research question.

\bibliographystyle{plain}
\bibliography{scholarai-paper-sample}
\end{document}
